\section{Conclusion}
\label{sec:conclusion}

We have argued, and shown on real EEG, that a decodable subject/domain signal is not by itself evidence
of a harmful functional shortcut. Measurable leakage magnitude does not certify functional reliance ---
task-head alignment is closer to reliance, though not a validated estimator. Subject signal is erasable,
but erasure strength does not certify a target benefit: no eraser achieves a proven gain across 40
held-out-target cells, and aggressive erasure can harm the target. A direct branch-local audit shows the
spatial branch is the strongest subject-leakage and load-bearing candidate, yet removing its subject
subspace \emph{hurts} the target --- so FSR refuses to call it a harmful shortcut; it is task-entangled.
And on injected positive controls, repair has a sharp scope: a controlled first-moment deterministic
offset is repairable by a deployable target-unlabeled mean alignment (narrowly, construction-matched),
whereas a controlled second-moment stochastic perturbation is not repaired even by an oracle-directed
covariance-shrinkage, and the natural signal is refused outright.

\begin{center}
\emph{Measurable is not reliance. Erasable is not beneficial. Natural leakage can be task-entangled, so
FSR refuses blind repair; and where repair holds, it holds only within a first-moment deterministic
scope.}
\end{center}

The contribution is a verification framework, not a new domain-generalization method: the FSR ladder
makes explicit what measurement licenses what control conclusion, and it draws the boundary of where
repair is --- and is not --- possible. FSR does not merely say ``do not repair''; it says \emph{when} to
refuse a blind repair (natural spatial leakage is task-entangled) and \emph{when} a repair holds only
within a mechanistic scope (first-moment deterministic offsets, not second-moment stochastic
perturbations, not natural or learned reliance). The claim ledger (Table~\ref{tab:claims}) and the
route-inclusion ledger (Appendix~\ref{app:routes}, Table~\ref{tab:routes}) record every claim's status,
so the boundary between what the evidence supports and what it does not is auditable.

The final picture is not that EEG subject leakage is harmless. It is that a shortcut claim requires an
information contract: leakage must be shown to be task-coupled, target-harmful, transferable, and
recoverable under the available observations. In our natural and controlled EEG audits --- and in the two
head-only learned-reliance probes --- those conditions \emph{split} rather than collapse into a single
leakage score.

\begin{table}[t]
\centering\small
\begin{tabular}{llp{7.6cm}}
\toprule
ID & Status & Claim \\
\midrule
C1 & READY\_WITH\_CAVEAT & Measured leakage magnitude does not certify reliance. \\
C2 & READY\_WITH\_CAVEAT & Task-head alignment is closer to reliance than raw leakage in frozen CIGL R3. \\
C3 & READY & Subject signal is erasable. \\
C4 & READY & Erasure strength does not certify target benefit. \\
C5 & READY\_WITH\_CAVEAT & Random-k falsifies nonspecific NLL movement. \\
C6 & READY & Spatial branch is load-bearing. \\
C7 & READY & Branch-local leakage/reliance is missing. \\
C8 & READY & CMI-control remains closed. \\
C9 & SUPPORT\_ONLY & TTA-Control is positive but non-CMI. \\
C10 & READY & FSR is an audit framework, not a new DG method. \\
\bottomrule
\end{tabular}

\caption{Claim ledger. Statuses are derived from the frozen result artifacts, not hand-set.}
\label{tab:claims}
\end{table}
